\chapter{Numbness}

\section*{Definition}
Numbness is reduced or absent sensation in part of the body. It may involve touch, pain, temperature, vibration, or position sense. Numbness can be temporary, localized to a compressed nerve, or a sign of disease affecting the brain, spinal cord, nerve roots, or peripheral nerves.

\section*{When to See a Doctor}
Call emergency services for sudden numbness of the face, arm, or leg, especially on one side, or numbness with weakness, speech difficulty, vision loss, severe headache, confusion, or loss of balance. Emergency assessment is also required for numbness after major injury, saddle numbness with bladder or bowel dysfunction, or a cold, pale, painful limb.

\section*{How It Develops}
Normal sensation depends on receptors, peripheral nerves, the spinal cord, and brain pathways. Pressure, inflammation, impaired blood flow, metabolic injury, or damage to any part of this system can reduce sensory signals. The distribution and speed of onset help distinguish a local nerve problem from a central neurologic emergency.

\section*{Causes}
\begin{itemize}
  \item Temporary nerve compression from posture, entrapment, or injury
  \item Diabetes, vitamin B12 deficiency, alcohol-related neuropathy, thyroid disease, or kidney disease
  \item Radiculopathy from a disk disorder, spinal stenosis, or shingles
  \item Stroke, transient ischemic attack, multiple sclerosis, seizure, or spinal cord compression
  \item Medication toxicity, autoimmune neuropathy, infection, or reduced circulation
\end{itemize}

\section*{Related Symptoms}
\begin{itemize}
  \item Tingling, burning, electric-shock sensations, pain, or loss of coordination
  \item Weakness, muscle wasting, gait difficulty, or impaired reflexes
  \item Back or neck pain, urinary changes, rash, or symptoms of diabetes
\end{itemize}

\section*{Medical Evaluation}
\begin{itemize}
  \item Examination maps sensation and tests strength, reflexes, coordination, gait, pulses, and cranial nerves.
  \item Blood tests may include glucose or A1c, vitamin B12, thyroid function, blood count, electrolytes, and targeted tests for inflammation or infection.
  \item Nerve-conduction studies, electromyography, MRI, or vascular imaging are selected according to the pattern and suspected location.
\end{itemize}

\section*{Treatment Options}
Treatment addresses the cause: relieving compression, correcting deficiencies, managing diabetes, treating infection or inflammation, providing stroke care, or treating spinal disease. Neuropathic pain may require clinician-directed medicines and physical or occupational therapy.

\section*{Self-Care and Home Management}
Protect numb areas from burns, cuts, pressure, and extreme temperatures. Change position regularly and avoid driving or operating machinery if sensation or strength is impaired. Do not dismiss progressive, recurrent, or one-sided numbness; arrange clinical assessment.

\section*{References}
\begin{thebibliography}{99}
\bibitem{numbness:ref1} England JD, et al. Practice parameter: evaluation of distal symmetric polyneuropathy. \textit{Neurology}. 2009;72:177--184.
\bibitem{numbness:ref2} Powers WJ, et al. 2019 Guidelines for the early management of acute ischemic stroke. \textit{Stroke}. 2019;50:e344--e418.
\bibitem{numbness:ref3} American Diabetes Association. Standards of Care in Diabetes: Microvascular Complications and Foot Care. 2025.
\end{thebibliography}
